\clearpage
\section{Validation}

The validation strategy is based on focused JUnit tests that exercise the
distributed protocol rather than only the local helper methods.
\begin{itemize}
    \item \texttt{DistributedCriticalSectionTest} validates concurrent
    bootstrap, mutual exclusion, crash recovery, repeated acquire/release
    cycles, invalid releases, and independence between distinct critical
    section names.
    \item \texttt{BrokerBootstrapLockTest} validates the temporary guard queue,
    including exclusive ownership and timeout behaviour.
    \item \texttt{TokenQueueManagerTest} validates the broker-side constraints
    of the token queue and the fail-fast behaviour on incompatible queue
    declarations.
\end{itemize}

This coverage is appropriate for the scope of the project because it targets
the exact risks of the design: duplicate bootstrap, token loss, inconsistent
release, and interference between unrelated critical sections. In other words,
the tests do not only check functional correctness, but also the safety
properties that justify the architecture itself.
